\begin{helper}
Let $G$ be a finite graph, let $\mu$ be an activation--priority independent-set
sample as in \noderef{RandomIndependentSetSampling}, and assume $\Delta>0$.
Fix three distinct pairwise nonadjacent vertices $a,b,c$, put
$Q=\{a,b,c\}$, and let
\[
\mathcal O=\{(a,b,c),(a,c,b),(b,a,c),(b,c,a),(c,a,b),(c,b,a)\}.
\]
For an order $t=(p,q,r)\in\mathcal O$, define the full-graph ordered-chamber
weight
\[
\begin{aligned}
\Phi_t
&=\frac{1}{\Delta^3}
\int_0^\gamma\int_z^\gamma\int_y^\gamma
\left(1-\frac{x}{\Delta}\right)^{|\{w\in V(G)\setminus\{p,q,r\}:w\sim p\}|} \\
&\qquad\qquad\cdot
\left(1-\frac{y}{\Delta}\right)^{
|\{w\in V(G)\setminus\{p,q,r\}:w\not\sim p,\ w\sim q\}|}
\left(1-\frac{z}{\Delta}\right)^{
|\{w\in V(G)\setminus\{p,q,r\}:w\not\sim p,\ w\not\sim q,\ w\sim r\}|}
\,dx\,dy\,dz .
\end{aligned}
\]
Then the survival probability of the fixed triple is bounded, as an extended
nonnegative real inequality, by the sum of the six full-graph ordered-chamber
weights:
\[
\operatorname{ofReal}\Pr[Q\subseteq I]\le
\operatorname{ofReal}\left(\sum_{t\in\mathcal O}\Phi_t\right).
\]
\end{helper}
\begin{proof}
Let $\nu$ be the activation--priority law supplied by
\noderef{RandomIndependentSetSampling}.  The push-forward identity
\noderef{SamplingFiniteSetPushForwardEvent} rewrites
\[
\operatorname{ofReal}\Pr[Q\subseteq I]
\]
as the $\nu$-measure of the event that every vertex of $Q$ belongs to the
output independent set.  Applying
\noderef{SamplingFiniteSetActivationPartition} partitions this event according
to the activation atom $A$.  Thus the left hand side is the sum, over
activation sets $A\supseteq Q$, of the fixed-activation event
\[
A\text{ is the activation set and every }q\in Q
\text{ beats every activated neighbor in }A .
\]

Fix such an activation set $A$.  Since $Q\subseteq A$, the three vertices
$a,b,c$ all lie in $A$; by hypothesis they are distinct and pairwise
nonadjacent.  The fixed-activation assembly helper
\noderef{SamplingTripleFixedActivationSixChamberAssembly} applies to this
atom.  It bounds the measure of the fixed-activation survival event by
\[
\sum_{t\in\mathcal O}\nu[A]\operatorname{ofReal}(W_t(A)),
\]
where $W_t(A)$ is the ordered chamber integral with the three atomwise
exponents
\[
|\{w\in A\setminus\{p,q,r\}:w\sim p\}|,\quad
|\{w\in A\setminus\{p,q,r\}:w\not\sim p,\ w\sim q\}|,
\]
and
\[
|\{w\in A\setminus\{p,q,r\}:w\not\sim p,\ w\not\sim q,\ w\sim r\}|.
\]
The same helper also states $W_t(A)\ge0$ for every order $t$.

Summing the fixed-activation inequalities over all $A\supseteq Q$ gives a
finite sum of extended nonnegative real terms.  For a fixed order $t$, the
atom measure in \noderef{RandomIndependentSetSampling} has the Bernoulli
form
\[
\nu[A]=\operatorname{ofReal}\left(
\left(\frac{\gamma}{\Delta}\right)^{|A|}
\left(1-\frac{\gamma}{\Delta}\right)^{|V(G)|-|A|}
\right).
\]
The Bernoulli atom weight is nonnegative, and the preceding paragraph gives
the nonnegativity of $W_t(A)$.  Therefore
\noderef{SamplingActivationENNRealWeightedSumConversion} converts, for each
fixed order $t$, the activation-atom sum of
$\nu[A]\operatorname{ofReal}(W_t(A))$ into the single extended nonnegative
real number obtained from the corresponding real Bernoulli-weighted activation
sum.

Finally apply \noderef{SamplingTripleActivationChamberSum} to that fixed
order $t=(p,q,r)$.  Its statement is exactly the activation-summation identity:
the real Bernoulli-weighted sum over activation atoms of the ordered chamber
integral $W_t(A)$ equals the full-graph ordered-chamber integral $\Phi_t$.
This use is valid for all six orders because each order is a permutation of
the distinct vertices $a,b,c$, and the pairwise nonadjacency hypothesis is not
needed by the summation identity beyond the sampling law itself.

Substituting these six identities into the summed fixed-activation bound gives
\[
\operatorname{ofReal}\Pr[Q\subseteq I]
\le
\operatorname{ofReal}\left(\sum_{t\in\mathcal O}\Phi_t\right),
\]
which is the claimed activation-summed six-order bound.  No paired chamber
normalization, infinite-chamber estimate, or final pair-sum reduction is used
in this helper.
\end{proof}
